\documentclass{beamer}
\title{Welcome to Investments!}
\begin{document}
\maketitle

\section{Pricing vs. Forecasting}

\begin{frame}{Asset Prices}
	\begin{itemize}
		\item An \underline{asset} is anything that produces something in the future
		\item Examples: stock, bond, tractor, milk cow
		\item A \underline{security} is a claim to the output of an asset
		\item The study and practice of investments is about \underline{pricing forecasts}, not \underline{forecasting prices}! 
		\item Example: Caterpillar (CAT) was paying quarterly dividends of \$0.86 per share. Will it continue to pay \$0.86 per share? What is the fair price of a claim to Caterpillar’s dividends (a share of stock)? Is the current market price the fair price?
		\item I will often refer to the study of investments as asset pricing
	\end{itemize}
\end{frame}

\section{Trading vs. Investing}

\begin{frame}{Two Hats}
	\begin{itemize}
		\item Think about wearing one of two hats:  \underline{investor} or \underline{speculator}
		\item Investors want to passively earn a return on their investments
		\item They trade so that they can share risk with other investors
		\item Speculators want to exploit small mispricings (hard to find)
		\item They trade so that they can buy from (sell to) other speculators or investors with worse information
	\end{itemize}
\end{frame}

\begin{frame}{Speculator Hat: \textit{where do trading profits come from?}}
	\begin{itemize}
		\item Trading (as typically understood) is a zero-sum game
		\item Ann \underline{knows} that stock $X$ is worth \$20
		\item Bob \underline{thinks} that $X$ is worth \$15
		\item Ann buys $X$ from Bob at \$15 and makes \$5
		\item At the end of the day, Mary at the factory was the only one who actually did any work
		\item Conclusion: people can make money trading
		\item However, acquiring information requires a lot of time and money, and even the best fail (e.g. LTCM, Great Recession)
	\end{itemize}
\end{frame}

\begin{frame}{Investor Hat: \textit{does that mean trading is bad? No.}}
	\begin{center}
		\renewcommand{\arraystretch}{1.5}
		\begin{tabular}{cc|cc}
			\multicolumn{2}{c}{\textbf{Bet A}} & \multicolumn{2}{c}{\textbf{Bet B}} \\
			Amount & Probability & Amount & Probability \\ \hline
			-\$10,000 & 0.001 & -\$10 & 1. \\
			\$0 & 0.999 & & \\ \hline
			\multicolumn{2}{c}{Average = -\$10} & \multicolumn{2}{c}{Average = -\$10}
		\end{tabular}
	\end{center}
\begin{center}
Which do you prefer?
\end{center}
\end{frame}

\begin{frame}{Investor Hat: \textit{does that mean trading is bad? No.}}
	\begin{itemize}
		\item Trading allows people to \underline{share risk}
		\item 100\% of Ann’s portfolio is in Exxon Mobil (XOM)
		\item 100\% of Bob’s portfolio is in Apple (AAPL)
		\item Let Ann and Bob trade so that they each have portfolios with 50\% in XOM and 50\% in AAPL
		\item It turns out that Ann and Bob reduce risk without sacrificing return. We’ll work out the math on this one later.
	\end{itemize}
\end{frame}

\section{Summary}

\begin{frame}{Summary}
	\begin{itemize}
		\item An asset is anything that produces something in the future
		\item A security is a claim to the output of an asset
		\item Think about two types of people: investors and speculators
		\item Investors want to passively earn a return on their investments
		\item They trade so that they can share risk with other investors
		\item Speculators want to exploit small mispricings (hard to find)
		\item They trade so that they can buy from (sell to) other speculators or investors with worse information
	\end{itemize}
\end{frame}
\end{document}
